% !TeX root = ../../Book3.tex
% 纲要标题：### 15.3 Koszul 正合性
% 纲要位置：工作记录/计划纲要.md，第 2462 行。
% 纲要细目（写作范围）：
% * 正则序列给出的自由消解
% * Noether 局部环上有限生成模的反向判据
% * Koszul 同调测量非正则性；零生成元、重复方程与冗余关系的正式算例
% * 逐元加入的短正合列、同调长正合列及一般与局部假设的对照

\section{Koszul 正合性}
\label{b3:sec:1503}

增加一个方程后，新的 Koszul 复形在每个次数包含原复形的相邻两项。两部分之间的微分由新增元素的乘法给出，因此新同调同时记录旧同调的商与乘法核。由此可以证明正则序列的正次同调消失；反过来，从同调消失恢复逐次单射，则需要局部有限性。


% 纲要内容（仅作写作计划，尚非教材正文）：
%
% * 正则序列给出的自由消解
% * Noether 局部环上有限生成模的反向判据
% * Koszul 同调测量非正则性
% * 逐元加入的短正合列、同调长正合列及一般与局部假设的对照
%

\subsection{正则序列给出的自由消解}
\label{b3:subsec:1503:01}

设 \(R\) 为交换环，\(M\) 为 \(R\) 模，\(r\ge1\)，\(x_1,\ldots,x_{r-1},y\in R\)，记 \(C_\bullet=K_\bullet(x_1,\ldots,x_{r-1};M)\)。增加 \(y\) 后，不含新增基向量的外积仍构成 \(C_i\)；含有它的外积则写成 \(e_y\wedge b\)，其中 \(b\in C_{i-1}\)。所以新复形包含一份原复形和一份平移后的原复形。把新增的基向量放在外积左端，得到
\[
K_i(x_1,\ldots,x_{r-1},y;M)\cong C_i\oplus C_{i-1},
\qquad \mathrm{d}(a,b)=(\mathrm{d}a+yb,-\mathrm{d}b).
\]
删去新增因子时产生 \(yb\)，而保留它时，旧微分前多一个负号；乘法 \(y\) 就这样连接两份复形。这里的链移位满足 \(C[1]_i=C_{i-1}\)、\(\mathrm{d}_{C[1]}=-\mathrm{d}_C\)。第一项的包含和第二项的投影都是链映射。

\begin{lemma}\label{b3:lem:koszul-adjoin-exact}
设 \(R\) 为交换环，\(M\) 为 \(R\) 模，\(r\ge1\)，\(x_1,\ldots,x_{r-1},y\in R\)，记 \(C_\bullet=K_\bullet(x_1,\ldots,x_{r-1};M)\)、\(K_\bullet=K_\bullet(x_1,\ldots,x_{r-1},y;M)\)。把新增基向量放在外积左端，识别 \(K_i=C_i\oplus C_{i-1}\)，并取 \(C[1]_i=C_{i-1}\)、\(\mathrm{d}_{C[1]}=-\mathrm{d}_C\)。以 \(\iota(a)=(a,0)\)、\(\pi(a,b)=b\) 为映射，有短正合列及其诱导的同调长正合列：
\[
\begin{gathered}
0\longrightarrow C_\bullet\xrightarrow{\iota}K_\bullet
\xrightarrow{\pi}C[1]_\bullet\longrightarrow0
\\
\Downarrow
\\
\cdots\longrightarrow H_i(C)\xrightarrow{y}H_i(C)
\longrightarrow H_i(K)
\longrightarrow H_{i-1}(C)\xrightarrow{y}H_{i-1}(C)\longrightarrow\cdots.
\end{gathered}
\]
\end{lemma}
\begin{proof}
对第二项投影得到的圈 \(b\in C_{i-1}\)，提升为 \((0,b)\)，其微分是 \((yb,0)\)。因此连接映射将 \([b]\) 送到 \([yb]\)，恰好是乘法 \(y\)。其余正合性也可从这一个公式核对：若圈 \((a,b)\) 投影为边界 \(b=\mathrm{d}c\)，则加上边界 \(\mathrm{d}(0,c)=(yc,-\mathrm{d}c)\) 后得到 \((a+yc,0)\)，来自 \(C\)；若 \(C\) 中圈 \(a\) 在 \(K\) 中是边界，则 \(a=\mathrm{d}u+yv\)、\(\mathrm{d}v=0\)，故其同调类来自乘法 \(y\)；最后，若 \(yb=-\mathrm{d}a\)，则 \((a,b)\) 是把 \([b]\) 提升到 \(K\) 的圈。这验证了长列中循环出现的三类位置。
\end{proof}

这条长列说明，同次旧同调模掉 \(y\) 的像后成为新同调的子模，而相邻低一次旧同调被 \(y\) 零化的部分成为其商。因此加入一元既可能消去旧同调，也可能从低一次同调的乘法核中产生新的同调。

\begin{theorem}\label{b3:thm:koszul-regular-resolution}
设 \(R\) 为交换环，\(M\ne0\) 为 \(R\) 模，\(r\ge0\) 为整数，\(x=(x_1,\ldots,x_r)\) 为 \(M\) 正则序列，\(I=(x_1,\ldots,x_r)\)。则
\[
H_i(x;M)=0\quad(i>0),\qquad H_0(x;M)=M/IM.
\]
增广复形 \(K_\bullet(x;M)\rightarrow M/IM\rightarrow0\) 正合，其次数 \(i\) 项为 \(M^{\binom ri}\)，一般不必自由。当 \(M=R\) 时，这是一条由有限秩自由模组成的 \(R/I\) 的自由消解。
\end{theorem}
\begin{proof}
当 \(r=0\) 时，复形只有次数零的 \(M\)，结论直接成立。一元情形已由乘法核算出。对长度归纳，令 \(C\) 如上。归纳假设给出 \(H_i(C)=0\) 对 \(i>0\)，而 \(H_0(C)=M/(x_1,\ldots,x_{r-1})M\)。上述正合列使新复形在 \(i\ge2\) 的同调为零，剩下的低次部分是
\[
0\longrightarrow H_1(K)\longrightarrow H_0(C)\xrightarrow{y}H_0(C)
\longrightarrow H_0(K)\longrightarrow0.
\]
归纳把剩余问题归结为 \(y\) 在这个逐次商上的乘法核。正则性使该核为零，次数零的余核则为 \(M/IM\)。增广映射是自然取商；Koszul 的次数 \(i\) 项为 \(\bigwedge^iR^r\otimes_RM\cong M^{\binom ri}\)，在 \(M=R\) 时各项自由。
\end{proof}

自由消解是由自由模组成的链复形，增广到被消解模以后正合。未增广的 Koszul 复形在次数零仍保留 \(H_0=M/IM\)，这里说它正合是指正次数同调消失；接上这个商与自然增广以后，才在所有位置正合。

\subsection{Noether 局部环上有限生成模的反向判据}
\label{b3:subsec:1503:02}

\begin{theorem}[Koszul 正合性判据]\label{b3:thm:koszul-acyclic}
设 \((R,\mathfrak m)\) 为交换 Noether 局部环，\(M\ne0\) 为有限生成的 \(R\) 模，\(r\ge0\) 为整数，\(x=(x_1,\ldots,x_r)\) 为 \(\mathfrak m\) 中的元素序列。以下条件等价：
\begin{equivalences}
\item \(x_1,\ldots,x_r\) 为 \(M\) 正则序列；
\item \(H_i(x;M)=0\) 对所有 \(i>0\) 成立；
\item \(H_1(x;M)=0\)。
\end{equivalences}
\end{theorem}
\begin{proof}
第一条件蕴含第二条件由\cref{b3:thm:koszul-regular-resolution}得到，第二条件显然蕴含第三条件。由第三条件对 \(r\) 归纳。\(r=0\) 时，空序列在非零模上正则；\(r=1\) 时，\(H_1(x_1;M)=0\) 就是乘法单射，而 \(x_1\in\mathfrak m\) 及 Nakayama 引理保证 \(M/x_1M\ne0\)。设 \(r>1\)，记 \(C=K_\bullet(x_1,\ldots,x_{r-1};M)\)、\(y=x_r\)。\cref{b3:lem:koszul-adjoin-exact}给出
\[
H_1(C)\xrightarrow{y}H_1(C)\longrightarrow0
\longrightarrow H_0(C)\xrightarrow{y}H_0(C).
\]
这同时给出 \(yH_1(C)=H_1(C)\) 及 \(y:H_0(C)\to H_0(C)\) 单射。每个 \(C_i\) 都是 \(M\) 的有限直和，因而是有限生成的 Noether 模；圈作为 \(C_1\) 的子模有限生成，边界作为 \(C_2\) 的同态像也有限生成，所以它们的商 \(H_1(C)\) 有限生成。对局部环 \(R\)、理想 \((y)\subseteq\mathfrak m=J(R)\) 及这个同调模使用\cref{b3:thm:nakayama}，得到 \(H_1(C)=0\)。归纳使前 \(r-1\) 项正则，而同时得到的单射正是 \(y\) 在相应逐次商上的非零因子条件。又 \(I=(x_1,\ldots,x_r)\subseteq\mathfrak m\)；若末端商 \(M/IM\) 为零，对有限生成模 \(M\) 再用 Nakayama 就会得到 \(M=0\)，矛盾。因此全部序列正则。
\end{proof}

这个证明只需次数一的消失；它通过有限性迫使较短序列的次数一同调也消失。若元素成为单位，乘法满射不能再由 Nakayama 推出模为零，这正是局部条件不能省略之处。

\ProseParagraphBreak
正则性中的末端非零条件和反向判据中的有限性条件承担不同的作用。\cref{b3:tab:koszul-hypotheses}把一般情形与本节的局部情形并列，便于区分它们。

\begin{center}
\begin{minipage}{\textwidth}
\makeatletter
\def\@captype{table}
\makeatother
\centering
\caption[正则序列与 Koszul 同调的假设和结论]{正则序列与 Koszul 同调的假设和结论。设 \(R\) 为交换环，\(M\ne0\) 为 \(R\) 模，\(x=(x_1,\ldots,x_r)\in R^r\)，\(I=(x_1,\ldots,x_r)\)。局部栏额外假设 \((R,\mathfrak m)\) 为 Noether 局部环、\(M\) 有限生成且全部 \(x_i\in\mathfrak m\)。局部化行假设原序列已经正则。}
\label{b3:tab:koszul-hypotheses}
\begin{tabularx}{\textwidth}{@{}>{\raggedright\arraybackslash}p{0.23\textwidth}>{\raggedright\arraybackslash}p{0.35\textwidth}>{\raggedright\arraybackslash}X@{}}
\toprule
条件或操作 & 一般环与模 & 上述 Noether 局部情形\\
\midrule
正则序列的条件 & 每一步乘法单射，且 \(M/IM\ne0\) & 每一步乘法单射；末端非零由 Nakayama 引理保证\\
正则推出 \(H_{i>0}=0\) & 成立 & 成立\\
\(H_{i>0}=0\) 推出正则 & 一般不成立 & 成立\\
\(H_1=0\) 推出正则 & 一般不成立 & 成立\\
任意重排保持正则 & 一般不成立 & 成立：同调不变性加本节的局部反向判据\\
在素理想 \(\mathfrak p\) 处局部化 & 保持逐次单射；还须 \((M/IM)_{\mathfrak p}\ne0\) & 在 \(\mathfrak p\in\operatorname{Supp}M\cap V(I)\) 处保持正则\\
\bottomrule
\end{tabularx}
\end{minipage}
\end{center}

\subsection{Koszul 同调与非正则性}
\label{b3:subsec:1503:03}

乘法出现新的零化关系，或新方程重复已有方程，都可能产生一次同调。下面几组计算分别说明这些同调类从何而来。

\begin{proposition}\label{b3:prop:koszul-obstruction-examples}
设 \(k\) 为域，\(R=k[X,Y]\)。对序列 \(X^2,XY\)，有
\[
\begin{gathered}
H_0(X^2,XY;R)\cong R/(X^2,XY),\\
H_1(X^2,XY;R)\cong R/(X),\qquad H_2(X^2,XY;R)=0.
\end{gathered}
\]
对序列 \(X,Y,0\) 及 \(X,Y,X+Y\)，零次和一次同调均同构于 \(k\)，二次以上同调为零。

若 \(R\) 为任意整环，\(z\in R\) 为非零非单位，则
\[
H_0(z,z;R)\cong R/(z),\qquad
H_1(z,z;R)\cong R/(z),\qquad H_2(z,z;R)=0.
\]
\end{proposition}
\begin{proof}
在 \(R=k[X,Y]\) 中，若 \((a,b)\in R^2\) 是 \(X^2,XY\) 的一次圈，则
\[
X^2a+XYb=0\quad\Longleftrightarrow\quad Xa+Yb=0.
\]
因 \(X\) 为素元且不整除 \(Y\)，等式迫使 \(b=Xc\)，再由整环中的消去得到 \(a=-Yc\)，且 \(c\) 唯一。因此圈模为 \(R(-Y,X)\)，边界模为 \(R(-XY,X^2)=XR(-Y,X)\)，得到所列一次同调。顶端微分 \(c\mapsto(-XYc,X^2c)\) 单射，故二次同调为零；一次微分的像为 \((X^2,XY)\)，给出零次同调。

在首次商 \(R/(X^2)\) 中，\(\bar X\ne0\) 而 \(XY\cdot\bar X=0\)。同样由 \(X,Y\) 互素，乘法 \(XY\) 的核恰为 \((X)/(X^2)\)；相邻两项的投影将圈 \(c(-Y,X)\) 的同调类送到 \(c\bar X\)，正把一次同调识别为这个乘法核。其支集为 \(V(X)\)，此处也恰好等于 \(V(X^2,XY)\)：在这条线的每个素点局部化后仍有非零一次同调，障碍并非只在原点。

序列 \(X,Y\) 在 \(R\) 上正则，商为 \(k\)，所以前一定理给出 \(C_\bullet=K_\bullet(X,Y;R)\) 的 \(H_0\cong k\) 及正次数同调消失。加入零元素后，微分公式中连接两份复形的乘法为零，故 \(K_\bullet(X,Y,0;R)=C_\bullet\oplus C[1]_\bullet\)。其同调为 \(H_i(C)\oplus H_{i-1}(C)\)，得到零次和一次各一个 \(k\)，其余为零。而
\[
\begin{pmatrix}1&0&0\\0&1&0\\-1&-1&1\end{pmatrix}
\begin{pmatrix}X\\Y\\X+Y\end{pmatrix}
=\begin{pmatrix}X\\Y\\0\end{pmatrix}.
\]
左侧矩阵行列式为 \(1\)，由\cref{b3:prop:koszul-generator-change}，序列 \(X,Y,X+Y\) 与 \(X,Y,0\) 具有同构的 Koszul 复形和相同的同调。

在任意整环 \(R\) 中，对非零 \(z\)，一次圈的条件 \(z(a+b)=0\) 等价于 \(a=-b\)，所以圈模为 \(R(-1,1)\)。边界模为 \(R(-z,z)=zR(-1,1)\)，其商是 \(R/(z)\)。顶端微分 \(c\mapsto(-zc,zc)\) 单射，一次微分的像为 \((z)\)，得到其余两项同调。\(z\) 非单位保证 \(R/(z)\ne0\)；第二次加入同一个方程不改变零次商，却留下 \((-1,1)\) 所生成的非零一次同调。
\end{proof}
